\section{Implementation boundaries and evidence}\label{sec:C20:implementation}
An inspectable harness exposes action parsing, context construction, tool dispatch, state persistence, stopping, and error recovery as separable modules. Code location alone is insufficient: a module can delegate essential behavior to a proprietary endpoint. The implementation record should identify the local revision, external service version where available, configuration, and which semantics remain opaque. This is the minimum needed to connect an algorithm to an identifiable execution path.

The OpenAI Agents SDK documentation describes an application-owned deployment path, whereas a managed agent service supplies a different runtime boundary~\cite{openai2026agents}. A thin SDK wrapper does not transfer ownership of application storage, tool semantics, or approvals to the provider. Conversely, a hosted runtime may maintain session state without exposing the underlying scheduler. These are role differences, not evidence that one architecture is universally more reliable. Chapter~\ref{ch:C22} compares their consequences without importing product priority rankings.

Harness\textsuperscript{2} illustrates adaptation of execution and improvement harnesses using behavioral comparisons and experience reuse~\cite{corpus261014224}. Without ground-truth task verifiers, such feedback remains vulnerable to shared model biases and execution variance. More editor sessions can improve the search over changes while also increasing cost and selection opportunities. A runtime supporting this mechanism needs candidate isolation, a record of rejected changes, and a rollback boundary that includes prompts, tools, and memory interpretation.

The current Harness\textsuperscript{2} source distinguishes logical stage accounting from all invocation costs. Its \path{core/ledger.py} checks rollout and improvement-stage identities $R=4k$ and $P=3k$ for recursion budget $k$, rejects duplicate stage labels, and separately records retry attempts. Selection and reflection are outside this ledger's stated counting boundary~\cite{cumb2026harness2code}. These are source-level count identities, not measurements of runtime or total expense. A complete cost comparison must join those stages with selection, reflection and failed invocations. The ledger is an adaptation-budget record, not a transactional journal proving remote-effect recovery.

An implementation audit should begin with one complete request, not an abstract class diagram. Trace the request from input serialization to terminal artifact, identify every mutable dependency, and enumerate where errors are caught or suppressed. A broad exception handler that returns a plausible textual tool result can turn infrastructure failure into misleading model evidence. A successful deserialization test will not expose that defect; an injected dependency outage will.

The strongest verification combines local contract checks with end-to-end behavioral experiments. Contract tests can establish that duplicate operation identifiers are rejected or that schema changes invalidate cached decisions. They cannot establish that a proposed patch satisfies an underspecified user task. Behavioral tests can estimate task reliability but may miss rare transactional windows. The two forms of evidence answer different questions and should remain separately reported.

The practical implication is a narrow but consequential design principle: authoritative execution state should be machine-checkable, recoverable, and distinct from model narration. The model remains essential for interpreting uncertain tasks and choosing useful actions. The runtime supplies the boundaries within which those choices can be evaluated and acted upon. Reliability emerges from their interaction; neither a fluent agent nor a durable database independently establishes it.
